We also prove that the Harsanyi dividend {\small$w_{\mathcal{S}}$} satisfies seven desirable axioms, including the \textit{efficiency, linearity, dummy, symmetry, anonymity, recursive} and \textit{interaction distribution} axioms, which demonstrates its trustworthiness.

(1) \textit{Efficiency axiom}. The output score of a model can be decomposed into effects of different causal patterns, \emph{i.e.} {\small $v(\boldsymbol{x})=\sum_{\mathcal{S}\subseteq\mathcal{N}}w_{\mathcal{S}}$}.

(2) \textit{Linearity axiom}. If we merge the output scores of the two models $t(\cdot)$ and $u(\cdot)$ into the output of model $v(\cdot)$, \emph{i.e.} {\small $\forall \mathcal{S}\subseteq\mathcal{N},~ v(\boldsymbol{x}_{\mathcal{S}})=t(\boldsymbol{x}_{\mathcal{S}})+u(\boldsymbol{x}_{\mathcal{S}})$}, the corresponding causal effects {\small $w^t_{\mathcal{S}}$} and {\small $w^u_{\mathcal{S}}$} can also be merged as {\small$\forall \mathcal{S}\subseteq \mathcal{N}, w^v_{\mathcal{S}}=w^t_{\mathcal{S}}+w^u_{\mathcal{S}}$}.

(3) \textit{Dummy axiom}. If a variable {\small $i\in\mathcal{N}$} is a dummy variable, \emph{i.e.} 
{\small $\forall \mathcal{S}\subseteq \mathcal{N}\backslash\{i\}, v(\boldsymbol{x}_{\mathcal{S}\cup\{i\}})=v(\boldsymbol{x}_{\mathcal{S}})+v(\boldsymbol{x}_{\{i\}})$}, it has no causal effect with other variables, {\small $\forall \mathcal{S}\subseteq \mathcal{N}\backslash\{i\}, w_{\mathcal{S}\cup\{i\}}=0$}.

(4) \textit{Symmetry axiom}. If the input variables {\small $i,j\in \mathcal{N}$} cooperate with other variables in the same manner, {\small $\forall \mathcal{S}\subseteq \mathcal{N}\backslash\{i,j\}, v(\boldsymbol{x}_{\mathcal{S}\cup\{i\}})=v(\boldsymbol{x}_{\mathcal{S}\cup\{j\}})$}, then they have the same causal effects with other variables, {\small $\forall S\subseteq \mathcal{N}\backslash\{i,j\}, w_{\mathcal{S}\cup\{i\}}=w_{\mathcal{S}\cup\{j\}}$}.

(5) \textit{Anonymity axiom}. For any permutations $\pi$ on {\small $\mathcal{N}$}, we have {\small $\forall \mathcal{S}
\!\subseteq\! \mathcal{N}, w^v_{\mathcal{S}}\!=\!w^{\pi v}_{\pi \mathcal{S}}$}, where {\small $\pi \mathcal{S} \!\triangleq\!\{\pi(i)|i\!\in\!\! \mathcal{S}\}$}, and the new model {\small $\pi v$} is defined by {\small $(\pi v)(\boldsymbol{x}_{\pi\mathcal{S}})\!=\!v(\boldsymbol{x}_{\mathcal{S}})$}.
This indicates that causal effects are not changed by the permutation.

(6) \textit{Recursive axiom}. The causal effects can be computed recursively.
For {\small $i\in \mathcal{N}$} and {\small $\mathcal{S}\subseteq \mathcal{N}\backslash\{i\}$}, the causal effect of the pattern {\small $\mathcal{S}\cup\{i\}$} is equal to the causal effect of {\small $\mathcal{S}$} in the presence of $i$ minus the causal effect of $\mathcal{S}$ in the absence of $i$, \emph{i.e.} {\small $\forall \mathcal{S}\!\subseteq\! \mathcal{N}\!\setminus\!\{i\}, w_{\mathcal{S}\cup\{i\}} = w_{\mathcal{S}|i~\text{present}} - w_{\mathcal{S}}$}. {\small $w_{\mathcal{S}|i~\text{present}}$} denotes the causal effect when the variable $i$ is always present as a constant context, \emph{i.e.} {\small $
w_{\mathcal{S}|i~\text{present}}=\sum_{\mathcal{S}'\subseteq \mathcal{S}} (-1)^{|\mathcal{S}|-|\mathcal{S}'|}\cdot v(\boldsymbol{x}_{\mathcal{S}'\cup\{i\}})$}.
